% figures/w05-ilos-idea.svg — 같은 조류에 LOS 와 ILOS 가 각각 무엇을 하는가.
%
%   bash _tools/tikz2svg.sh figures/src/w05-ilos-idea.tex figures/w05-ilos-idea.svg
%
% 기하를 눈대중으로 찍지 않는다. 강의 §5-5 의 측정에서 V_c = 0.3 m/s 일 때 배가
% 붙드는 선수각은 -22.8 deg 이고 Delta = 5 m 이므로 정상상태 오프셋은
%
%     y_e^ss = Delta tan(beta) = 5 * tan(22.8 deg) = 2.10 m
%
% 이고 강의 표의 측정치는 2.107 m 이다 (§5-5 의 실험 5-5a). 그림의 각과 길이는
% 그 값 그대로 쓴다 — 0.005 m 의 차이는 인쇄된 선수각을 소수점 한 자리로
% 반올림한 데서 온다.
%
% ── 부호를 맞춰 둔다 ────────────────────────────────────────────────────────
% 침로가 경로와 나란하면 chi = pi_p 이고 chi = psi + beta 이므로 psi = pi_p - beta:
% 선수는 경로 방향에서 **반시계로** beta 만큼 돌아간다. NED 에서 양의 각은
% 시계방향이므로, 화면에서 경로를 오른쪽(동쪽)으로 그리면 선수는 **위로** 든다.
% 한편 y_e^{\,p} = +Delta tan(beta) 는 경로 방향의 **우현** 쪽, 곧 화면 아래쪽이다.
% 그래서 LOS 선체는 경로 아래에 있고 선수는 위를 향한다 — 조류를 거슬러.
% 조류는 우현으로 밀므로 화면에서 아래로 그린다.
%
% ALOS 판은 이 주차가 가르치지 않는다 (W05_simulink/_previous_version/ 에 있다).
% 세 법칙을 함께 그린 옛 그림은 figures/src/w05-ilos-alos-idea.tex 에 남겨 둔다.
%
% 출처: Fossen (2021) Handbook §12.3 (LOS),
%       Borhaug, Pavlov & Pettersen (2008) CDC, pp. 4984-4991 (ILOS).
\documentclass[border=5pt]{standalone}
% gnc-style.tex — 이 볼트의 TikZ 그림이 공유하는 스타일.
%
% 저널 그림 규약을 스타일 하나로 굳혀 둔다. 그림마다 선 굵기나 화살촉을
% 다시 정하지 않는다 — 그것이 그림이 제각각으로 보이는 원인이다.
%
% 쓰는 법:  % gnc-style.tex — 이 볼트의 TikZ 그림이 공유하는 스타일.
%
% 저널 그림 규약을 스타일 하나로 굳혀 둔다. 그림마다 선 굵기나 화살촉을
% 다시 정하지 않는다 — 그것이 그림이 제각각으로 보이는 원인이다.
%
% 쓰는 법:  % gnc-style.tex — 이 볼트의 TikZ 그림이 공유하는 스타일.
%
% 저널 그림 규약을 스타일 하나로 굳혀 둔다. 그림마다 선 굵기나 화살촉을
% 다시 정하지 않는다 — 그것이 그림이 제각각으로 보이는 원인이다.
%
% 쓰는 법:  \input{gnc-style}  (figures/src/ 안에서)

\usepackage{tikz}
\usepackage{amsmath}
\usepackage{xcolor}
\usetikzlibrary{arrows.meta, positioning, calc, fit, backgrounds, decorations.pathmorphing}

% ---- 색 --------------------------------------------------------------------
% 색은 강조 하나에만 쓴다. 나머지는 검정.
\definecolor{ink}{HTML}{1A1A1A}
\definecolor{accent}{HTML}{7C3AED}
\definecolor{warn}{HTML}{B91C1C}
\definecolor{soft}{HTML}{666666}

% 기하 그림(좌표계·유도)에서만 쓰는 넷. 블록선도는 위의 흑백 + accent 로 충분하다.
\definecolor{pathc}{HTML}{0D9488}   % 계획 경로 · 항적
\definecolor{errc}{HTML}{C2410C}    % 오차 (cross-track 등)
\definecolor{alng}{HTML}{B45309}    % along-track
\definecolor{flow}{HTML}{0369A1}    % 조류 · 유속

% ---- 선과 화살촉 -----------------------------------------------------------
\tikzset{
  % 신호선. 화살촉은 작고 뾰족한 것 하나뿐이다.
  sig/.style   = {ink, line width=0.5pt, -{Stealth[length=4pt, width=3pt]}},
  wire/.style  = {ink, line width=0.5pt},
  % 강조 경로 — 파선으로 구분한다. 흑백 인쇄에서도 살아야 하기 때문이다.
  asig/.style  = {accent, line width=0.5pt, densely dashed,
                  -{Stealth[length=4pt, width=3pt]}},
  awire/.style = {accent, line width=0.5pt, densely dashed},
  % 블록. 흰 바탕, 직각, 검은 테두리.
  blk/.style   = {draw=ink, line width=0.55pt, fill=white,
                  minimum width=17mm, minimum height=8mm, inner sep=2pt, font=\small},
  % 합산점
  sum/.style   = {draw=ink, line width=0.55pt, fill=white, circle,
                  inner sep=0pt, minimum size=4.6mm, font=\scriptsize},
  asum/.style  = {draw=accent, line width=0.55pt, fill=white, circle,
                  inner sep=0pt, minimum size=4.6mm, font=\scriptsize},
  % 분기점
  dot/.style   = {ink, fill=ink, circle, inner sep=0pt, minimum size=1.5mm},
  % 신호 이름 — 선 위에 작게
  sname/.style = {font=\small, inner sep=1.5pt},
  % 그림 안의 짧은 설명. 그림 안의 글은 최소로 둔다
  note/.style  = {font=\scriptsize, soft, inner sep=1.5pt},
  anote/.style = {font=\scriptsize, accent, inner sep=1.5pt},
  % 축
  axis/.style  = {ink, line width=0.5pt},
  tick/.style  = {font=\scriptsize, soft},
}

% 캡션 한 줄 — 그림 아래 가는 선과 함께
\newcommand{\gnccaption}[2]{%
  \node[anchor=north west, text width=#1, align=left, font=\scriptsize, ink]
        at (current bounding box.south west) {\vspace{2pt}#2};}
  (figures/src/ 안에서)

\usepackage{tikz}
\usepackage{amsmath}
\usepackage{xcolor}
\usetikzlibrary{arrows.meta, positioning, calc, fit, backgrounds, decorations.pathmorphing}

% ---- 색 --------------------------------------------------------------------
% 색은 강조 하나에만 쓴다. 나머지는 검정.
\definecolor{ink}{HTML}{1A1A1A}
\definecolor{accent}{HTML}{7C3AED}
\definecolor{warn}{HTML}{B91C1C}
\definecolor{soft}{HTML}{666666}

% 기하 그림(좌표계·유도)에서만 쓰는 넷. 블록선도는 위의 흑백 + accent 로 충분하다.
\definecolor{pathc}{HTML}{0D9488}   % 계획 경로 · 항적
\definecolor{errc}{HTML}{C2410C}    % 오차 (cross-track 등)
\definecolor{alng}{HTML}{B45309}    % along-track
\definecolor{flow}{HTML}{0369A1}    % 조류 · 유속

% ---- 선과 화살촉 -----------------------------------------------------------
\tikzset{
  % 신호선. 화살촉은 작고 뾰족한 것 하나뿐이다.
  sig/.style   = {ink, line width=0.5pt, -{Stealth[length=4pt, width=3pt]}},
  wire/.style  = {ink, line width=0.5pt},
  % 강조 경로 — 파선으로 구분한다. 흑백 인쇄에서도 살아야 하기 때문이다.
  asig/.style  = {accent, line width=0.5pt, densely dashed,
                  -{Stealth[length=4pt, width=3pt]}},
  awire/.style = {accent, line width=0.5pt, densely dashed},
  % 블록. 흰 바탕, 직각, 검은 테두리.
  blk/.style   = {draw=ink, line width=0.55pt, fill=white,
                  minimum width=17mm, minimum height=8mm, inner sep=2pt, font=\small},
  % 합산점
  sum/.style   = {draw=ink, line width=0.55pt, fill=white, circle,
                  inner sep=0pt, minimum size=4.6mm, font=\scriptsize},
  asum/.style  = {draw=accent, line width=0.55pt, fill=white, circle,
                  inner sep=0pt, minimum size=4.6mm, font=\scriptsize},
  % 분기점
  dot/.style   = {ink, fill=ink, circle, inner sep=0pt, minimum size=1.5mm},
  % 신호 이름 — 선 위에 작게
  sname/.style = {font=\small, inner sep=1.5pt},
  % 그림 안의 짧은 설명. 그림 안의 글은 최소로 둔다
  note/.style  = {font=\scriptsize, soft, inner sep=1.5pt},
  anote/.style = {font=\scriptsize, accent, inner sep=1.5pt},
  % 축
  axis/.style  = {ink, line width=0.5pt},
  tick/.style  = {font=\scriptsize, soft},
}

% 캡션 한 줄 — 그림 아래 가는 선과 함께
\newcommand{\gnccaption}[2]{%
  \node[anchor=north west, text width=#1, align=left, font=\scriptsize, ink]
        at (current bounding box.south west) {\vspace{2pt}#2};}
  (figures/src/ 안에서)

\usepackage{tikz}
\usepackage{amsmath}
\usepackage{xcolor}
\usetikzlibrary{arrows.meta, positioning, calc, fit, backgrounds, decorations.pathmorphing}

% ---- 색 --------------------------------------------------------------------
% 색은 강조 하나에만 쓴다. 나머지는 검정.
\definecolor{ink}{HTML}{1A1A1A}
\definecolor{accent}{HTML}{7C3AED}
\definecolor{warn}{HTML}{B91C1C}
\definecolor{soft}{HTML}{666666}

% 기하 그림(좌표계·유도)에서만 쓰는 넷. 블록선도는 위의 흑백 + accent 로 충분하다.
\definecolor{pathc}{HTML}{0D9488}   % 계획 경로 · 항적
\definecolor{errc}{HTML}{C2410C}    % 오차 (cross-track 등)
\definecolor{alng}{HTML}{B45309}    % along-track
\definecolor{flow}{HTML}{0369A1}    % 조류 · 유속

% ---- 선과 화살촉 -----------------------------------------------------------
\tikzset{
  % 신호선. 화살촉은 작고 뾰족한 것 하나뿐이다.
  sig/.style   = {ink, line width=0.5pt, -{Stealth[length=4pt, width=3pt]}},
  wire/.style  = {ink, line width=0.5pt},
  % 강조 경로 — 파선으로 구분한다. 흑백 인쇄에서도 살아야 하기 때문이다.
  asig/.style  = {accent, line width=0.5pt, densely dashed,
                  -{Stealth[length=4pt, width=3pt]}},
  awire/.style = {accent, line width=0.5pt, densely dashed},
  % 블록. 흰 바탕, 직각, 검은 테두리.
  blk/.style   = {draw=ink, line width=0.55pt, fill=white,
                  minimum width=17mm, minimum height=8mm, inner sep=2pt, font=\small},
  % 합산점
  sum/.style   = {draw=ink, line width=0.55pt, fill=white, circle,
                  inner sep=0pt, minimum size=4.6mm, font=\scriptsize},
  asum/.style  = {draw=accent, line width=0.55pt, fill=white, circle,
                  inner sep=0pt, minimum size=4.6mm, font=\scriptsize},
  % 분기점
  dot/.style   = {ink, fill=ink, circle, inner sep=0pt, minimum size=1.5mm},
  % 신호 이름 — 선 위에 작게
  sname/.style = {font=\small, inner sep=1.5pt},
  % 그림 안의 짧은 설명. 그림 안의 글은 최소로 둔다
  note/.style  = {font=\scriptsize, soft, inner sep=1.5pt},
  anote/.style = {font=\scriptsize, accent, inner sep=1.5pt},
  % 축
  axis/.style  = {ink, line width=0.5pt},
  tick/.style  = {font=\scriptsize, soft},
}

% 캡션 한 줄 — 그림 아래 가는 선과 함께
\newcommand{\gnccaption}[2]{%
  \node[anchor=north west, text width=#1, align=left, font=\scriptsize, ink]
        at (current bounding box.south west) {\vspace{2pt}#2};}


\def\BE{22.8}        % 크랩각 [deg] — 강의 §5-5 의 측정
\def\SC{6}           % 1 m = 6 mm
\def\OFF{12.6}       % Delta tan(beta) = 2.10 m -> 12.6 mm

\tikzset{
  pth/.style  = {pathc, line width=0.8pt, densely dashed},
  cur/.style  = {flow, line width=0.9pt, -{Stealth[length=4.2pt,width=3.0pt]}},
  bow/.style  = {ink, line width=0.9pt, -{Stealth[length=5pt,width=3.6pt]}},
  trk/.style  = {pathc, line width=1.0pt, -{Stealth[length=5pt,width=3.6pt]}},
}
% 선체 — 경로 방향에서 반시계로 beta 만큼 돌아간 채로 (x,y) 에 놓는다
\newcommand{\hull}[2]{%
  \begin{scope}[shift={(#1,#2)}, rotate=\BE]
    \fill[ink, opacity=0.88] (7.5,0) -- (-1.5,3.0) -- (-1.5,-3.0) -- cycle;
  \end{scope}}
\newcommand{\hd}[2]{\parbox[b]{86mm}{\raggedright
  {\small\bfseries\color{ink}#1}\\[1.5pt]{\scriptsize\color{soft}#2}}}
% 높이를 못 박는다 — \parbox[t] 만으로는 tabular 가 아래 깊이를 세지 않는다
\newcommand{\ft}[1]{\parbox[t][22mm][t]{86mm}{\raggedright\scriptsize\color{ink}#1}}

\begin{document}
\begin{tabular}{@{}l@{}}

\parbox{180mm}{\raggedright
  {\small\bfseries\color{ink}
   Both laws end up pointing the bow upstream by the same angle. They differ in
   what is left over.}\\[2pt]
  {\scriptsize\color{soft}
   Same vessel, same path, same current: $V_c = 0.3$\,m/s on the beam, so the
   vessel must hold $\beta = 22.8^\circ$ of crab and $\Delta\tan\beta = 2.10$\,m
   with $\Delta = 5$\,m. Experiment 5-5a measures the offset at $2.107$\,m.
   \textbf{The bow tilt is the same in both panels — only one of them is on the
   path.}}}\\[3.5mm]

\begin{tabular}{@{}c@{\hspace{6mm}}c@{}}
\hd{1\ \ LOS}{the lean is spent on the drift}
& \hd{2\ \ ILOS}{an error that is not there, so the real one can be zero}
\\[1.5mm]

% ===================== 1 : LOS =============================================
\begin{tikzpicture}[x=1mm, y=1mm, baseline={(0,0)}]
  \path (-2,-30) rectangle (86,26);
  \draw[pth] (0,0) -- (84,0);
  \node[note, pathc, anchor=south west] at (0,0.8) {the path};

  % 조류 — 경로의 우현(화면 아래)으로 민다
  \foreach \x in {8,26,44,62} {
    \draw[cur] (\x,20) -- (\x,13);}
  \node[note, flow, anchor=west] at (2,23) {the current, $V_c = 0.3$\,m/s};

  % 선체 : 경로 아래 y_e^{\,p} = +2.10 m
  \hull{24}{-\OFF}
  \draw[bow] (24,-\OFF) -- ++({26*cos(\BE)},{26*sin(\BE)});
  % 화살촉 **오른쪽**에 붙인다. south east 로 두면 이름표가 화살대 위에 얹힌다
  % (CLAUDE.md §4 규칙 6) / to the right of the arrowhead: anchored south east the
  % label lies on top of the arrow shaft
  \node[note, ink, anchor=west] at ({24+26*cos(\BE)+1},{-\OFF+26*sin(\BE)}) {bow, $\psi$};
  \draw[trk] (24,-\OFF) -- (58,-\OFF);
  \node[note, pathc, anchor=north] at (44,{-\OFF-1}) {track, $\chi$};
  \draw[accent, line width=0.6pt] ({24+17},{-\OFF}) arc (0:\BE:17);
  \node[anote, anchor=west] at ({24+18.5},{-\OFF+4}) {$\beta$};

  % 남는 오프셋
  \draw[errc, line width=0.7pt, {Stealth[length=3.6pt]}-{Stealth[length=3.6pt]}]
        (70,0) -- (70,-\OFF);
  \node[font=\scriptsize, errc, anchor=west] at (71,{-0.5*\OFF}) {$y_e^{\,p} = 2.10$\,m};
\end{tikzpicture}
&
% ===================== 2 : ILOS ============================================
\begin{tikzpicture}[x=1mm, y=1mm, baseline={(0,0)}]
  \path (-2,-30) rectangle (86,26);
  \draw[pth] (0,0) -- (84,0);
  \node[note, pathc, anchor=south west] at (0,0.8) {the path};

  \foreach \x in {8,26,44,62} {
    \draw[cur] (\x,20) -- (\x,13);}
  \node[note, flow, anchor=west] at (2,23) {the current, $V_c = 0.3$\,m/s};

  % 선체 : 경로 위, y_e^{\,p} = 0
  \hull{24}{0}
  \draw[bow] (24,0) -- ++({26*cos(\BE)},{26*sin(\BE)});
  \node[note, ink, anchor=south west] at ({24+26*cos(\BE)},{26*sin(\BE)}) {bow, $\psi$};
  \draw[trk] (24,0) -- (58,0);
  \node[note, pathc, anchor=north] at (44,-1) {track, $\chi$};
  \draw[accent, line width=0.6pt] (41,0) arc (0:\BE:17);
  \node[anote, anchor=west] at (42.5,4) {$\beta$};

  % 있지도 않은 오차
  \draw[accent, line width=0.7pt, densely dashed] (2,-\OFF) -- (70,-\OFF);
  \draw[accent, line width=0.6pt, {Stealth[length=3.6pt]}-{Stealth[length=3.6pt]}]
        (70,0) -- (70,-\OFF);
  \node[anote, anchor=east] at (69,{-0.5*\OFF}) {$\kappa\,y_{int} = 2.10$\,m};
  \node[anote, anchor=north west, align=left, text width=62mm] at (2,{-\OFF-1.5})
       {a \textbf{phantom} error the integrator built. There is no vessel here.};
\end{tikzpicture}
\\[1mm]
\ft{The bow is tilted upstream by $22.8^\circ$ — but that tilt \emph{is} the
    $\arctan(y_e^{\,p}/\Delta)$ term, and it is entirely used up cancelling the
    drift. Nothing is left to close $y_e^{\,p}$. The vessel runs \textbf{parallel}
    to the path, $2.107$\,m to the side of it, for ever (Experiment 5-5a).}
& \ft{The law reads $y_e^{\,p} + \kappa y_{int}$, not $y_e^{\,p}$. The integrator
      fills until its phantom error equals the old $2.10$\,m offset, and only then
      can the \textbf{real} error be zero — measured at $0.011$\,m with
      $\kappa = 0.3$ (Experiment 5-5b). It never learns what a current is; it only
      knows the error must vanish.}
\\
\end{tabular}\\[3.5mm]

\parbox{180mm}{\raggedright\scriptsize
  \textbf{The one idea to take away.} A current cannot be removed by steering
  harder. It must be \textbf{answered} by pointing the bow upstream. LOS points
  upstream but spends the whole lean on the drift, so an offset is left over;
  ILOS buys that lean from a second term and hands the arctan its original job
  back.\\[3pt]
  \textbf{And the anti-windup is in the law, not bolted on.} The integrator's own
  update is
  $\dot y_{int} = \Delta y_e^{\,p} / (\Delta^2 + (y_e^{\,p} + \kappa y_{int})^2)$,
  whose denominator grows with the error — so the state barely fills while the
  vessel is far off the path, and fills fastest just where it should, close to it.
  Weeks 2 to 4 had to add back-calculation for the same effect.\\[3pt]
  \textbf{Sources.} LOS — Fossen (2021), \emph{Handbook of Marine Craft
  Hydrodynamics and Motion Control}, 2nd ed., §12.3. ILOS — B\o rhaug, Pavlov
  and Pettersen (2008), \emph{Proc.\ 47th IEEE CDC}, pp.~4984--4991. The numbers
  in the panels are §5-5's measurements, not illustrations.}

\end{tabular}
\end{document}
